\subsection{Centralizers and free class-two source subgroups}
\label{subsec:cm-class-two-subgroups}

Retain the hypotheses and notation of
Theorem~\ref{thm:cm-measure-class-two}, and write
\[
 Q=Q_2^{\rm CM},\qquad Z=\operatorname{im}J_A,
 \qquad \pi:Q\longrightarrow A.
\]
The following statements use the proved integral map from the
actual source group to this quotient.

\begin{proposition}\label{prop:cm-class-two-centralizers}
The centre of \(Q\) is exactly \(Z\).  For every \(q\in Q\)
with \(a=\pi(q)\ne0\),
\[
 C_Q(q)=\pi^{-1}\bigl(A\cap\mathbb Q_p a\bigr).
\]
The intersection is taken inside \(A[1/p]\), is a free
\(\mathbb Z_p\)-module of rank one, and the displayed
centralizer is, noncanonically, \(Z\times\mathbb Z_p\).
In particular, every abelian subgroup of \(Q\) has image
of \(\mathbb Z_p\)-rank at most one in \(A\).
\end{proposition}

\begin{proof}
The commutator of two elements with abelian images \(a,b\)
is \(J_A(a\wedge b)\).  Proposition~\ref{prop:cm-decomposable-commutators}, together
with injectivity of \(\alpha[1/p]\), gives the exact criterion
\[
 J_A(a\wedge b)=0\quad\Longleftrightarrow\quad
 b\in\mathbb Q_pa\qquad(a\ne0).
\]
This proves the asserted centralizer formula.

Here is the integral topology of the line intersection.
Choose a continuous \(\mathbb Z_p\)-linear functional
\(\ell:A\to\mathbb Z_p\) with \(\ell(a)\ne0\).
Such a functional exists because the compact torsion-free
abelian pro-\(p\) group \(A\) is a product of copies of
\(\mathbb Z_p\).  If \(\lambda a\in A\), then
\(\lambda\ell(a)\in\mathbb Z_p\), which bounds the
denominator of \(\lambda\).  The allowed scalars form a
nonzero closed \(\mathbb Z_p\)-submodule of this bounded
rank-one lattice.  They are consequently \(p^e\mathbb Z_p\)
for some integer \(e\).  Thus \(A\cap\mathbb Q_pa\)
has a generator \(a_0\).  Choose \(q_0\in Q\) above
\(a_0\).  The continuous powers of \(q_0\) identify
its closed procyclic group with \(\mathbb Z_p\), and
\((z,t)\mapsto zq_0^t\) is an isomorphism from
\(Z\times\mathbb Z_p\) onto the displayed centralizer.

The module \(A[1/p]\) has dimension at least two: for
example, the moment-isolated elements constructed below
have distinct nonzero initial homogeneous terms.
Every nonzero \(a\) therefore has a \(\mathbb Q_p\)-independent
\(b\), so an element above \(a\) is not central.
Hence the centre is precisely \(Z\).  The assertion about
abelian subgroups follows from the same centralizer formula.
\end{proof}

\begin{corollary}\label{cor:cm-three-generator-class-two}
Let \(r\le3\), and let \(h_1,\ldots,h_r\in H\) have
\(\mathbb Q_p\)-linearly independent images in \(A[1/p]\).
For the closed subgroup \(U\) they generate, the map
\(U\to Q\) has kernel exactly \(\gamma_3U\), and
\(U/\gamma_3U\) is the free class-two pro-\(p\) group
on these \(r\) generators.
\end{corollary}

\begin{proof}
Let \(B\) be the \(\mathbb Z_p\)-span of their abelian
images.  Every bivector in the exterior square of a
vector space of dimension at most three over
\(\mathbb Q_p\) is decomposable.  The preceding
derivative argument thus makes
\(J_A:\bigwedge^2B\to Z\) injective.
Map the free pro-\(p\) group \(P_r\) to \(H\) using
the chosen elements.  The composite to \(Q\) kills
\(\gamma_3P_r\); its maps on the abelian quotient and
on \(\gamma_2P_r/\gamma_3P_r\simeq\bigwedge^2\mathbb Z_p^r\)
are injective.  Its kernel is therefore exactly
\(\gamma_3P_r\).  The image of that subgroup in \(U\)
is \(\gamma_3U\), which proves both conclusions.
\end{proof}

\begin{proposition}\label{prop:cm-countable-class-two}
There is a closed subgroup \(U\subset H\), generated
by a sequence tending to one, for which
\[
 U\cap N_2^{\rm CM}=\gamma_3U,
 \qquad U/\gamma_3U\simeq P/\gamma_3P,
\]
where \(P\) is the free pro-\(p\) group on a countably
infinite convergent sequence of generators.
Thus the complete free class-two quotient of this
actual source subgroup is detected.  Every finite
initial subsequence gives the corresponding conclusion
in finite rank.
\end{proposition}

\begin{proof}
For \(i\ge1\), put \(k_i=2^i\) and
\(d_i=2k_i-2\).  The character
\(\kappa_{(k_i,k_i)}\) is nonzero by
\cite[Theorem 2.8(2) and the Soul\'e-finiteness paragraph,
p.~9]{Ishii}; this is the ordinary \(k_i\)-th
Tate twist, with \(k_i\ge2\).

Let
\[
 \mathcal S_i=\{\boldsymbol m\in I:
                              |\boldsymbol m|\le2k_i\}.
\]
Put \(\Gamma_2=\operatorname{Gal}(L_\infty/L_1)\), and choose \(\gamma_i\in\Gamma_2\) for which the finitely
many characters \(\chi_{\boldsymbol m}(\gamma_i)\),
\(\boldsymbol m\in\mathcal S_i\), are pairwise distinct.
The open rank-two principal-unit image of the two CM
Tate characters, established by the ray-class morphism
in \cite[pp.~4--5]{Ishii}, permits this choice:
the forbidden loci after taking logarithms are a finite
union of proper linear subspaces.  Put
\[
 P_i(X)=\prod_{\substack{\boldsymbol m\in\mathcal S_i\\
                         \boldsymbol m\ne(k_i,k_i)}}
                 (X-\chi_{\boldsymbol m}(\gamma_i)).
\]
Choose \(b_i^0\in A\) with
\(\kappa_{(k_i,k_i)}(b_i^0)\ne0\), and set
\(b_i=P_i(\gamma_i)b_i^0\).  Equivariance shows that
all its moments in \(\mathcal S_i\), except the desired
one, vanish, and that the desired one remains nonzero.

Choose a lift \(\widetilde h_i\in H\) of \(b_i\), and
put \(h_i=\widetilde h_i^{p^i}\), with abelian image
\(a_i=p^ib_i\).  This sequence tends to one: in every
finite quotient of \(H\), a sufficiently large power
\(p^i\) kills every element.  The exponential transform
\(f_i\) of \(\alpha(a_i)\) has the form
\[
 f_i(U_1,U_2)=c_i(U_1U_2)^{k_i-1}
                    +\text{terms of total degree greater than }d_i,
 \qquad c_i\in\mathbb Q_p^\times.
\]
In particular, the \(d_i\) are strictly increasing.

Let \(P\) be free on the convergent sequence
\(x_i\to1\).  Its abelianization is
\(\prod_{i\ge1}\mathbb Z_p x_i\), and the integral
class-two Hall calculation gives
\[
 \gamma_2P/\gamma_3P
       =\prod_{i<j}\mathbb Z_p[x_i,x_j].
\]
These assertions follow by taking inverse limits of
the corresponding finite-generator, finite-quotient
Hall presentations.  The continuous map \(x_i\mapsto h_i\)
induces an injective abelian map: in a nonzero completed
linear combination of the \(f_i\), the least index
with nonzero coefficient supplies a unique nonzero
initial degree.

The initial homogeneous term of \(J(f_i,f_j)\) is
\[
 c_ic_j J\bigl((U_1U_2)^{k_i-1},
              (U_1U_2)^{k_j-1}\bigr),
\]
of degree
\[
 d_i+d_j=2^{i+1}+2^{j+1}-4.
\]
It is nonzero by the derivative computation above:
the two displayed monomials have nonconstant ratio.
These degrees are pairwise distinct for \(i<j\), by
uniqueness of binary expansion.  A nonzero completed
alternating combination therefore again has a unique
least initial degree, which cannot cancel.
For each total degree there are only finitely many
pairs to consider, so this argument applies to the
completed product, not only to finite combinations.
Consequently the composite
\(P\to H\to Q\) induces an injection on both its
abelian quotient and its class-two central quotient.
Its kernel is exactly \(\gamma_3P\).

Let \(U\) be the image of \(P\) in \(H\).  It is compact
and hence closed.  The image of \(\gamma_3P\) is
\(\gamma_3U\), by the commutator definition and
compactness.  The exact kernel just calculated gives
\(U\cap N_2^{\rm CM}=\gamma_3U\) and the asserted
isomorphism.  No injectivity assertion for the full
map \(P\to H\) is needed for these conclusions.
\end{proof}

The subgroups in Proposition~\ref{prop:cm-countable-class-two}
depend on the moment-isolating choices.  The surrounding
quotient \(Q_2^{\rm CM}\) remains the canonical arithmetic
quotient of Theorem~\ref{thm:cm-measure-class-two}.
